Table~\ref{tab:results} reports measurements from the fixed protocol. The
shared model reaches $-22.24\pm0.20$~dB on held-out IID trajectories, compared
with $-14.92\pm0.35$~dB for the approximately parameter-matched MLP and
$-21.19\pm0.30$~dB for the standalone recurrent model. Thus recurrence is a
strong alternative to the extra slow tier in this example. The shared model
has 28,031 parameters, whereas the standalone recurrent control has 5,398.
The latter should remain a serious baseline when a slow semantic/context tier
is not otherwise required by the application.

Under the joint traffic, coupling, and mobility shift, the shared predictor's
NMSE degrades to $-9.58\pm0.41$~dB. The standalone recurrent model reaches
$-9.43\pm0.37$~dB. This small difference does not establish a robust
out-of-distribution benefit from the nested hierarchy. Persistence is also
competitive with the smaller MLP under this shift, illustrating that training
accuracy alone is an inadequate model-selection argument for deployment.

Sharing changes the IID mean only modestly relative to the unshared model,
whose three-seed variability is substantially larger. Removing modulation
reduces IID accuracy, but its OOD mean is close to the full model. Most
importantly, the correct bridge comparison---unshared joint versus unshared
detached---does not establish a consistent advantage for cross-scale gradients
across IID and OOD conditions. The experiment therefore supports the feasibility
of the training graph and APIs, rather than a general superiority claim for
each proposed coupling.

The paired action-effect NMSE for the shared model is
$-11.75\pm0.27$~dB, below the action-ignorant persistence control's 0~dB.
This provides a direct check that the learned predictor uses candidate actions.
It does not by itself establish an intervention effect in a real radio system;
the comparison is against the explicitly intervenable synthetic simulator.

Table~\ref{tab:intent} also reveals substantial parsing error despite the small
controlled vocabulary. The shared model's mean TV increases from 0.118 for
held-out wording to 0.195 for the unseen three-objective combination. An
additional exact lexicon compiler achieves zero TV on all of these restricted
grammar splits, as recorded in \texttt{intent\_baseline.json}. That control
requires no neural training and makes clear that the present language task
does not justify using a pretrained language model. The learned head is a
demonstration of differentiable intent compilation whose robustness needs
improvement and evaluation on a less artificial corpus.

The shared model's finite-candidate IID regret is
$0.0055\pm0.0022$ in normalized utility units. All selected candidates satisfy
the configured power budget because infeasible candidates are removed before
ranking. Small regret on this short candidate list should not be confused with
high long-run policy return or learned constraint satisfaction.

On the measured CPU host (Python 3.12.14, PyTorch 2.5.1 CPU, one PyTorch thread),
the shared model's mean per-run fast p50 is approximately 0.081~ms and its mean
per-run fast p99 is approximately 0.291~ms. The slow p50 is approximately
0.591~ms. These are measurements of a tiny eager-mode model, with the fast
packet already cached. They are neither a comparison with foundation-model
inference nor a worst-case deadline guarantee. Raw timing samples and
per-initialization quantiles are preserved, so the host-specific observation
can be checked without treating it as a deployed baseband result.
